
\section{Deriving the Model}
    
    \subsection{Model Ingredients}

        We now review several model classes that provide useful components for constructing a probabilistic emulator.

        \subsubsection{Partition Models}

            Partition models divide the covariate space into a collection of disjoint regions, each associated with a separate statistical model. 
            The simplest example is a regular grid in $\mathcal{E}$ where each bin contains a histogram or kernel density estimate of the observed outputs.

            A more flexible construction uses a Voronoi tessellation. 
            Given a set of seeds $\{\vec{s}_i\}_{i=1}^{N_d} \subset \mathcal{E}$, each point $\vec{x}$ is assigned to the nearest seed according to a distance metric:
            \begin{equation}
                \text{cell}(\vec{x}) =
                \argmin_i
                \left\{
                    \sqrt{(\vec{x}-\vec{s}_i)^\intercal G_i (\vec{x}-\vec{s}_i)}
                \right\}.
            \end{equation}

            Here $G_i$ is a positive-definite metric tensor associated with the $i^{\text{th}}$ cell. 
            In many practical implementations (\href{https://www.jstor.org/stable/pdf/48583762.pdf}{such as the paper we have been following}) this metric is restricted to a diagonal form
            \begin{equation}
                G_i = \mathrm{diag}(\vec{w}_i)
            \end{equation}
            with $\sum_j [w_i]_j = 1$, which preserves the usual geometric properties of a Voronoi tessellation such as linear cell boundaries and simply connected regions.

            Partition models are attractive because they allow different statistical behaviour in different regions of the parameter space. 
            However, they suffer from a major limitation: predictions are typically discontinuous across cell boundaries. 
            Capturing smoothly varying behaviour therefore requires a large number of partitions.

        \subsubsection{Mixture of Experts}

            An alternative approach is the \emph{mixture of experts} model, in which the conditional distribution is represented as a weighted sum of simpler component models:
            \begin{equation}
                p(\y|\vec{x}) =
                \sum_{i=1}^{N_e} w_i(\vec{x}) p_i(\y|\vec{x}),
            \end{equation}
            where $p_i$ are expert distributions and the weights satisfy
            \begin{equation}
                \sum_i w_i(\vec{x}) = 1.
            \end{equation}

            In many machine learning applications the weights are generated by a gating network, for example
            \begin{equation}
                w_i(\vec{x}) =
                \frac{\exp(\vec{k}_i^\intercal \vec{x}+b_i)}
                {\sum_j \exp(\vec{k}_j^\intercal \vec{x}+b_j)}.
            \end{equation}

            More complex architectures allow the covariates $\vec{x}$ to first be transformed through multiple hidden layers before the mixture weights are computed.

            Mixture models provide smooth interpolation between experts, but the gating mechanism is often purely algebraic and does not explicitly exploit the geometric structure of the covariate space. In short, it is a black box that `works' with little explainability as to why one expert was more heavily weighted than the other.

            \begin{analogy}
                This is termed a `mixture of experts' because it is helpful to think of the process as a set of nodes which are making predictions as to the answer, and then reaching a consensus based on how close the query point is to each of their `areas of expertise': the closer to an expert the point is, the more weight their voice is given to the final consensus.
            \end{analogy}
        \subsubsection{Kernel-Based Expert Interpolation}

            A simpler and more geometrically interpretable weighting scheme can be constructed using kernels. 
            Suppose each expert is associated with a location $\vec{c}_i$ in the covariate space. 
            We define weights using a radial basis function kernel:
            \begin{equation}
                w_i(\vec{x}) =
                \frac{K\!\left(d(\vec{x},\vec{c}_i \mid G)\right)}
                {\sum_{\ell} K\!\left(d(\vec{x},\vec{c}_\ell \mid G)\right)}.
            \end{equation}

            Here
            \begin{equation}
                d(\vec{x},\vec{y}\mid G)
                =
                \sqrt{(\vec{x}-\vec{y})^\intercal G (\vec{x}-\vec{y})}
            \end{equation}
            is a Mahalanobis distance induced by the metric $G$, and $K(d)$ is a radial kernel satisfying $K(0)=1$ and decaying with increasing distance.

            Under this construction the prediction at $\vec{x}$ is obtained by smoothly interpolating between nearby experts. 
            If $\vec{x}$ lies close to $\vec{c}_k$ then $w_k(\vec{x}) \approx 1$ and the prediction is dominated by the $k^{\text{th}}$ expert.

    \subsection{Department-Expert Models}

        A Department-Expert Model combines the geometric structure of partition models with the smooth interpolation properties of kernel mixtures.

        We assume that the model contains

        \begin{itemize}
            \item a set of Voronoi tessellation seeds $\{\vec{s}_k\}_{k=1}^{N_d}$, and
            \item a set of experts $\{\vec{c}_i,p_i\}_{i=1}^{N_e}$ located within the covariate space.
        \end{itemize}

        If the partitioning were applied in a \emph{hard} manner, predictions would only depend on experts located in the same Voronoi cell as the query point, giving the following form:

        \begin{align}
            p({y}|\vec{x}) &= \sum_{k=1}^{N_d} T_k(\vec{x}) \sum_{i=1}^{N_e} w^k_i(\vec{x}) p_i(\y),
            \\
            w^k_i(\vec{x}) &=
            \frac{T_k(\vec{c}_i) K(d(\vec{x},\vec{c}_i))}
            {\sum_j T_k(\vec{c}_j) K(d(\vec{x},\vec{c}_j))},
            \\
            T_k(\vec{a}) & = \delta_{k,\text{cell}(\vec{a})} =  \begin{cases} 1 & \text{cell}(\vec{a}) = k \\ 0 & \text{else} \end{cases}
        \end{align}

        This is identical to the pure mixture model, but with $T_k$ removing elements from the sum if the points do not belong to $k$ - and since the only non-zero contribution to the outer sum is from the $k$ to which $\x$ belongs, it ensures that $p(\y|\x)$ only gets contributions from experts in the same Voronoi cell.
        
        \begin{analogy}
            To continue the analogy of `experts' reaching a consensus, we term each of the cells generated by this partition as a \textit{department}. Each point in the covariate space -- and each expert -- is then assigned a `department', and only the experts in that department may contribute to the consensus.
        \end{analogy}

        This is a \textit{Strict Department of Expert} (SDE) model.


        However, SDE models suffer from discontinuities in the likelihood whenever $\vec{x}$ crosses a cell boundary. 
        These discontinuities make the likelihood surface non-differentiable and therefore unsuitable for gradient-based optimisation.

        To avoid this difficulty we introduce a \emph{flexible department} border. Rather than $T_k$ acting to assigning each point to a single cell, we define it as a probability that a point belongs to the department:
        \begin{align}
            T_k(\vec{x}) = P(\text{cell}(\vec{x}) = k | \vec{\vT}) &=
            \frac{K_2\!\left(d(\vec{x},\vec{s}_k)\right)}
            {\sum_j K_2\!\left(d(\vec{x},\vec{s}_j)\right)}.
        \end{align}

        Here $K_2$ is the `department kernel' a distinct radial basis function that we choose \textit{a priori}. When the kernel $K_2$ is sharply peaked the model approaches an SDE, while broader kernels produce smoother transitions between regions of parameter space. This defines a \textit{Flexible Department of Experts} (FDE) model.

        \begin{analogy}
            In the analogy: experts are no longer purely confined to one department, but rather spend varying amounts of time visting each department. The more time they spend in one department (corresponding to a higher $T_k$), the more weight their voice is given in a consensus. 
        \end{analogy}
        

    \subsection{A Flexible \& Anisotropic Department of Experts}

        In the basic FDE model formulation above, we assumed that the metric $G$ was constant across space, and thus that $d$ was similarly constant. 

        The final ingredient of our model -- the Flexible \& Anisotropic Department of Experts (\name{}) Model -- is to equip each department with a learned metric. Each seed $\vec{s}_k$ is associated with a metric tensor $G_k$, and the distance between two covariate vectors is evaluated locally as
        \begin{equation}
            d_k(\vec{x},\vec{y})
            =
            \sqrt{(\vec{x}-\vec{y})^\intercal G_k (\vec{x}-\vec{y})}.
        \end{equation}

        If $G_k$ is restricted to a diagonal form (as in  \href{https://www.jstor.org/stable/pdf/48583762.pdf}{Payne et al.}) this simply rescales the coordinate axes, assigning different relative importance to each covariate within the $k^{\text{th}}$ region. 

        Relaxing this restriction allows $G_k$ to contain off-diagonal terms, which encode correlations between covariates and therefore rotate and stretch the local coordinate system. Consequently each department does not merely partition the parameter space but instead defines a local \emph{correlation topology}. 

        \begin{analogy}
            To finalise the analogy of a FADE model - the varying metric means that each department is associated with a unique way of interpreting the covariate space and the correlations within it.

            In this analogy, each department has a unique interpretation of the same technical terms - the word `stress' might mean something different to a physicist ("force per unit area"), a biologist (a physiological response to a stimuli) and a linguist (the prominence given to a syllable in a word). When an expert is `working in' a department, they must use the terminology of that department: an expert close to the border of two departments will rapidly switch beetween terminologies -- whilst those ensconsed in their own department will use only their own jargon.
        \end{analogy}
        The full specification of a \name{} model is detailed in table \ref{T:fullModel}
        \begin{table*}
            \large
            \begin{align}
            p^\text{fade}({y}|\vec{x},\vec{\Delta},\vT,\vec{\eta}) &=
            \sum_{k=1}^{N_d} T_k(\x) \sum_{i=1}^{N_e}
            \,   w_i^k(\vec{x})\, p_i(\y),
            \\
            w_i^k(\vec{x}) &=
            \frac{T_k(\vec{c}_i) K_1(d_k(\vec{x},\vec{c}_i))}
            {\sum_{j=1}^{N_e} T_k(\vec{c}_j) K_1(d_k(\vec{x},\vec{c}_j))},
            \\
            T_k(\vec{x}) & =  \frac{K_2(d_k(\vec{x},\vec{s}_k))}
            {\sum_{j=1}^{N_d} K_2(d_j(\vec{x},\vec{s}_j))},
            \\
            d_k(\x,\y) & = \sqrt{(\x-\y)^\intercal G_k (\x -\y)}
        \end{align}
        \hrule
        
        \begin{flushleft}
        \small  Where:
        \begin{enumerate}
            \item $\y \in \mathcal{Y} \subset \mathbb{R}^q$ is the value to be emulated. The space $\mathcal{Y}$ is constrained to be a subset of $\mathbb{R}^q$ by the choice of $\mathcal{F}$ and the prior.
            \item $\vec{x} \in \mathcal{E} \subset \mathbb{R}^n$ is the covariate input vector
            \item $\vec{\Delta} = (N_d,N_e)$ are our discrete model parameters
            \begin{enumerate}
                \item $N_d$ is the number of Voronoi seeds used (due to the anisotropic metric, the space may be split into $n \geq N_d$ disjoint regions)
                \item $N_e$ is the number of experts populating the space
            \end{enumerate}
            \item $\vec{\eta}= (K_1,K_2,\mathcal{F})$ are the model hyperparameters - our final result will be conditioned on these values.
            \begin{enumerate}
                \item $K_1(d)$ is the \textit{expert kernel}, an isotropic distance kernel of some kind, used to determine how quickly the information is exchanged between experts in the same department
                \item $K_2(d)$ is the \textit{department kernel}, an isotropic distance kernel used to determine how easily experts can communicate between departments
                \item $\mathcal{F}$ is a family of probability distribution functions that each expert may choosen from.
            \end{enumerate}
            \item $\vT = \left(\left\{\vec{s}_k,\{\phi^k_{ij}\}\right\}_{k=1}^{N_{p}},\{\vec{c}_i, \vec{\xi}_i\}_{i=1}^{N_e}\right)$ specifies our partitioned-expert scheme
            \begin{enumerate}
                \item $\vec{s}_k \in \mathcal{E}$ are the coordinates seeds of the $k$\textsuperscript{th} Voronoi pseudo-partition
                % \item $L_k  \in \mathbb{R}^{n\times n}$ is the Cholesky decomposition of the metric, $G_k$ 
                \item $\{\phi^k_{ij}\}_{0 \leq i < n}^{0 \leq j \leq i}$ are a set of parameters ($n(n+1)/2$ of them) which are used to define a lower triangular matrix with positive diagonal elements:
                \begin{equation}
                    [L_k]_{ij} = \begin{cases} \exp(\phi^k_{ij}) & i = j \\ \phi^k_{ij} & i > j \\ 0 &\text{else}\end{cases}
                \end{equation}
                This then defines a unique metric, $G_k(\{\phi^k\}) = L_k L_k^\intercal$, which is guaranteed to be positive definite and symmetric. $\phi^k_{ij}$ is otherwise unconstrained
                \item $\vec{c}_i \in \mathbb{R}^n$ are the coordinates of the $i$\textsuperscript{th} expert.
                \item $\vec{\xi}_i$ is the parameter set which specifies $p_i$ in the family $\mathcal{F}$:
                \begin{equation}
                    p_i(\y) = \mathcal{F}(y; \vec{\xi}_i)
                \end{equation}
                In the case where $\vec{y} = y \in \R$ and $\mathcal{F}$ is a univariate Gaussian, $\vec{\xi}_i = (\mu_i, \sigma_i)$.
            \end{enumerate}
        \end{enumerate}
    \end{flushleft}
    \hrule
        \caption{The full specification of a \name{} model.}\label{T:fullModel}
    \end{table*}
 
